\chapter{Time integration}
\label{ch:integration}
\impl{src/integrators.jl}

Three integrators share one \code{predict!}\,/\,\code{evaluate!}\,/\,%
\code{correct!} protocol, so the nonlinear and linear machinery downstream is
written once.

\section{Quasi-static}
\label{sec:quasistatic}

Inertia is dropped from \eqref{eq:semidiscrete} and the load is applied in
pseudo-time increments:
\begin{equation}
  \Rvec(\Uvec) = \Fext(t_{n+1}) - \Fint(\Uvec) = \vect{0},
  \qquad \Kmat = \frac{\partial\Fint}{\partial\Uvec} .
\end{equation}
The operator is the bare tangent stiffness.  For a well-shaped mesh its
condition number grows like $(L/h)^2$ with the mesh refinement ratio, which is
what makes quasi-static problems the natural home for multigrid
(Chapter~\ref{ch:amg}) and the reason a diagonal preconditioner degrades so
visibly with problem size.

\section{Newmark-\texorpdfstring{$\beta$}{beta}}
\label{sec:newmark}

The Newmark family parameterizes the step by $\beta$ and $\gamma$:
\begin{align}
  \Uvec_{n+1} &= \Uvec_n + \Delta t\, \Vvec_n
                + \Delta t^2\left[\left(\tfrac12 - \beta\right)\Avec_n + \beta \Avec_{n+1}\right], \\
  \Vvec_{n+1} &= \Vvec_n + \Delta t\left[(1-\gamma)\Avec_n + \gamma \Avec_{n+1}\right].
\end{align}
Carina implements this in predictor--corrector form.  The predictor collects
everything known at $t_n$,
\begin{equation}
  \Uvec^{\mathrm{pred}} = \Uvec_n + \Delta t\, \Vvec_n
                          + \Delta t^2\left(\tfrac12 - \beta\right)\Avec_n,
  \qquad
  \Vvec^{\mathrm{pred}} = \Vvec_n + \Delta t (1-\gamma)\Avec_n,
\end{equation}
after which the acceleration follows from the displacement alone:
\begin{equation}
  \Avec_{n+1} = c_M\left(\Uvec_{n+1} - \Uvec^{\mathrm{pred}}\right),
  \qquad
  \boxed{\ c_M = \frac{1}{\beta \Delta t^2}\ }
  \label{eq:cM}
\end{equation}
and the corrector restores the velocity, $\Vvec_{n+1} = \Vvec^{\mathrm{pred}} +
\gamma\Delta t\,\Avec_{n+1}$.  Displacement is thus the only unknown, and the
Newton system for it has the \emph{effective operator}
\begin{equation}
  \Keff = \Kmat + c_M \Mmat .
  \label{eq:keff}
\end{equation}

Equation~\eqref{eq:cM} governs almost everything about how implicit dynamics
behaves numerically, and it is worth being explicit about why.  As $\Delta t$
shrinks, $c_M \sim \Delta t^{-2}$ grows without bound and $\Keff$ becomes
mass-dominated.  The mass matrix is well conditioned --- it is a Gram matrix of
shape functions, independent of the deformation --- so $\Keff$ inherits that
conditioning.  Three consequences recur throughout this manual:

\begin{itemize}
  \item A diagonal (Jacobi) preconditioner is unusually effective at small
        $\Delta t$, because $\diag(c_M\Mmat)$ captures the dominant term.
  \item Multigrid, whose value is mesh-independent iteration counts on a
        stiffness-dominated operator, has little left to remove; measurements
        confirm it does not repay its per-application cost in this regime
        (Section~\ref{sec:amg-when}).
  \item A low-rank quasi-Newton model of $\Keff^{-1}$ is viable at small
        $\Delta t$ and not at all in the quasi-static limit
        (Section~\ref{sec:lbfgs}).
\end{itemize}

\paragraph{Constrained degrees of freedom.}  The predictor is engineered so
that \eqref{eq:cM} reproduces the prescribed acceleration exactly at
constrained nodes: $\Uvec^{\mathrm{pred}}_{\mathrm{p}} = \Uvec_{\mathrm{p}} -
\ddot{g}(t_{n+1})/c_M$, whence $\Avec_{\mathrm{p}} = \ddot{g}(t_{n+1})$.  This
keeps the inertial cross-term $\Mmat_{\mathrm{fp}}\ddot{g}$ in the residual
without special-casing it, and it is why the integrator carries full-length
$\Uvec$, $\Vvec$, $\Avec$ arrays rather than free-DOF slices.

\paragraph{Stability.}  $\beta = \tfrac14$, $\gamma = \tfrac12$ (average
acceleration) is unconditionally stable and second-order accurate, and is the
default.  Numerical dissipation, when wanted, is available through the
HHT-$\alpha$ generalization the integrator carries.

\section{Central difference}
\label{sec:central-difference}

The explicit integrator advances
\begin{equation}
  \Uvec_{n+1} = \Uvec_n + \Delta t \Vvec_n + \tfrac12 \Delta t^2 \Avec_n,
  \qquad
  \Avec_{n+1} = \left(\Mmat^{\mathrm{L}}\right)^{-1}\Rvec(\Uvec_{n+1}),
\end{equation}
followed by $\Vvec_{n+1} = \Vvec_n + \Delta t\left[(1-\gamma)\Avec_n + \gamma
\Avec_{n+1}\right]$.  With a lumped mass the acceleration update is a
component-wise division: no linear solve, no nonlinear solve, no tangent.
Nothing in Chapters~\ref{ch:newton}--\ref{ch:amg} applies, which is why the
input parser ignores the \code{solver} block entirely for this integrator.

\paragraph{Critical time step.}  Stability is conditional on the CFL limit
\begin{equation}
  \Delta t \le \mathrm{CFL} \cdot \min_e \frac{\ell_e}{c_e},
  \qquad c_e = \sqrt{\frac{\lambda + 2\mu}{\rho_0}},
\end{equation}
with $\ell_e$ a characteristic element length and $c_e$ the dilatational wave
speed.  Carina recomputes this bound periodically during the run rather than
once at setup, because $\ell_e$ shrinks as elements compress.

\section{Adaptive time stepping}
\label{sec:adaptive}

Both implicit integrators shrink $\Delta t$ by a decrease factor when the
nonlinear solve fails and grow it by an increase factor, up to a ceiling, when
it succeeds.  A failure is any of: the status tree reporting
\textsc{Failed}~(Section~\ref{sec:status}), a non-finite residual, or a domain
error raised inside a constitutive model.

The mechanism matters most for $J_2$ plasticity, where it is not merely an
efficiency device but a correctness one.  Because the return map is path
dependent, a step large enough to traverse the yield surface in one increment
integrates the wrong path --- and converges perfectly well to the wrong answer.
Step-size control plus the line search of Section~\ref{sec:linesearch} is what
keeps the increments small enough for the return map to track the true path.
